\section{Dados e métodos}

\subsection{Conjuntos de dados}

Os dados estão organizados em quatro arquivos, cada um com três colunas: o índice de
tempo discreto $n$, a entrada $u[n]$ e a saída $y[n]$. A Tabela~\ref{tab:dados}
resume os conjuntos e o uso de cada um.

\begin{table}[H]
\centering
\caption{Conjuntos de dados utilizados no trabalho.}
\label{tab:dados}
\begin{tabular}{lcl}
\toprule
Arquivo & Amostras & Uso \\
\midrule
\texttt{dados\_treino\_branco.csv}    & 5000 & Treino com entrada branca \\
\texttt{dados\_validacao\_branco.csv} & 3000 & Validação do modelo \\
\texttt{dados\_treino\_colorido.csv}  & 5000 & Treino com entrada colorida \\
\texttt{dados\_teste\_anomalia.csv}   & 5000 & Teste de detecção de anomalia \\
\bottomrule
\end{tabular}
\end{table}

\subsection{Métodos de estimação}

As estatísticas de primeira e segunda ordem foram estimadas diretamente das amostras.
A média e a variância seguem os estimadores usuais. A autocorrelação foi calculada
com o estimador enviesado a partir dos sinais centrados, e normalizada pelo valor no
atraso zero para facilitar a leitura. Para a correlação cruzada adotou-se a definição
$R_{yu}[k] = \mathrm{E}\{y[n]\,u[n-k]\}$, de modo que, para entrada branca, a parte
causal aparece nos atrasos positivos.

No domínio da frequência, a densidade espectral de potência foi estimada pelo método
de Welch, assim como o espectro cruzado e a coerência. Não há informação sobre a
frequência de amostragem, então as frequências são apresentadas de forma normalizada,
em ciclos por amostra, no intervalo de $0$ a $0{,}5$.

A resposta em frequência foi estimada por dois caminhos. O primeiro usa apenas as
densidades espectrais e fornece o módulo,
\begin{equation}
|\hat{H}_1(e^{j\omega})| = \sqrt{\frac{\hat{S}_{yy}(e^{j\omega})}{\hat{S}_{uu}(e^{j\omega})}} .
\end{equation}
O segundo usa o espectro cruzado e fornece também a fase,
\begin{equation}
\hat{H}_2(e^{j\omega}) = \frac{\hat{S}_{yu}(e^{j\omega})}{\hat{S}_{uu}(e^{j\omega})} .
\end{equation}
A coerência, por sua vez, é definida por
\begin{equation}
\gamma^2_{uy}(e^{j\omega}) = \frac{|\hat{S}_{yu}(e^{j\omega})|^2}{\hat{S}_{uu}(e^{j\omega})\,\hat{S}_{yy}(e^{j\omega})} ,
\end{equation}
e indica, em cada frequência, o quanto a saída está linearmente relacionada com a
entrada.

\subsection{Modelo identificado}

Optou-se por um modelo de resposta ao impulso finita (FIR),
$\hat{y}[n] = \sum_{k=0}^{M-1} h[k]\,u[n-k]$, com os coeficientes $h[k]$ estimados por
mínimos quadrados. Essa escolha se ajusta bem ao caso de entrada branca, no qual a
correlação cruzada já fornece uma estimativa direta da resposta ao impulso, e mantém
a interpretação simples, sem a necessidade de escolher ordens de partes
autorregressivas ou a dimensão de um espaço de estados.
